%==============================================================================
%  (iii)  Spectral proper orthogonal decomposition
%
%  Theory and the estimator actually used.  No results here.
%==============================================================================
\section{Spectral proper orthogonal decomposition}
\label{sec:spod}

%------------------------------------------------------------------------------
\subsection{Formulation}
\label{sec:spod:formulation}

Let $\vect{q}(\xvec,t)$ denote the fluctuation of the state vector about the
long-time mean, $\vect{q} = \vect{u} - \meanx{\vect{u}}$. For a flow that is
statistically stationary in time and homogeneous in the streamwise direction,
the two-point space--time correlation depends on $t$ and $x$ only through the
separations $\tau$ and $\Delta x$, so the correlation tensor is diagonalised by
a Fourier transform in both. Writing $\qhat(\kx,\cdot,\omega)$ for the
transformed fluctuation, the cross-spectral density tensor at each
wavenumber--frequency pair is
\begin{equation}
  \csd(\kx,\omega)
    = \avg{\, \qhat(\kx,\omega)\,\qhat^{*}(\kx,\omega) \,} ,
  \label{eq:csd}
\end{equation}
where $(\cdot)^{*}$ denotes the conjugate transpose and $\avg{\cdot}$ the
expectation over realisations.

SPOD modes are the eigenvectors of \cref{eq:csd} in the energy inner product
$\ip{\vect{q}_1}{\vect{q}_2} = \vect{q}_2^{*}\Wmat\vect{q}_1$, with $\Wmat$ the
positive-definite matrix of cell volumes:
\begin{equation}
  \csd(\kx,\omega)\,\Wmat\,\spodmode_j(\kx,\omega)
    = \spodev_j(\kx,\omega)\,\spodmode_j(\kx,\omega) ,
  \label{eq:spod-eig}
\end{equation}
with the eigenvalues ordered $\spodev_1 \ge \spodev_2 \ge \dots \ge 0$.
The modes are orthonormal, $\ip{\spodmode_i}{\spodmode_j} = \delta_{ij}$, each
$\spodev_j$ is the expected energy of the flow projected onto $\spodmode_j$, and
the expansion is optimal in the sense that no other orthonormal basis captures
more energy in the first $n$ modes at that $(\kx,\omega)$
\citep{Lumley1970,Towne2018}.

\TODO{Add the one sentence that makes the choice of SPOD --- rather than
space-only POD or DMD --- non-arbitrary for this flow: SPOD modes at a single
frequency evolve coherently in time and are the structures that a
statistically stationary flow actually sustains, whereas space-only POD mixes
frequencies into a single mode.}

%------------------------------------------------------------------------------
\subsection{Estimation from the simulation}
\label{sec:spod:estimator}

\TODO{The estimator, stated so that it is reproducible.  Welch's method
\citep{Welch1967}: $\nblk$ blocks of $\nfft$ snapshots with $50\,\%$ overlap and
a \TODO{Hamming} window; the resulting frequency resolution
$\Delta f = 1/(\nfft \Delta t_s)$ and bandwidth; the discrete FFT in $x$ over
the $N$ sampled planes, giving $\kx = 2\pi n/L_x$ for $n = 0,\dots,N/2$; and the
eigendecomposition of the $\nblk \times \nblk$ Gram matrix rather than of
$\csd$ itself, since $\nblk \ll$ the number of grid points.}

\begin{equation}
  \csd(\kx,\omega) \approx
    \frac{1}{\nblk}\sum_{m=1}^{\nblk}
      \qhat^{(m)}(\kx,\omega)\,\qhat^{(m)*}(\kx,\omega) ,
  \label{eq:welch}
\end{equation}

\TODO{Convergence: how many blocks were available, and the resulting uncertainty
on $\spodev_1$.  State honestly which modes are converged --- typically the
leading one --- and which are reported as indicative only.  Note that $\pm\kx$
are complex conjugates and so do not supply additional realisations.}

\TODO{Normalisation and presentation conventions: whether spectra are premultiplied
by $\omega$, how modes are phase-aligned before plotting, and the energy measure
used ($\spodev_1/\sum_j \spodev_j$ as the low-rank indicator).}
